\chapter{微分形式}\label{chap:differential-forms}
\begingroup
\renewcommand{\thefootnote}{}
\footnotetext{\textcopyright\ Springer International Publishing Switzerland 2014.
  L. Godinho and J. Natário, \emph{An Introduction to Riemannian Geometry},
  Universitext. DOI \href{https://doi.org/10.1007/978-3-319-08666-8_2}{10.1007/978-3-319-08666-8\_2}.}
\endgroup

本章讨论可微流形上的积分. 由于没有典范的局部坐标选取, 也就没有自然的体积概念, 因而只有在坐标变换下具有适当变换性质的对象才能积分. 这些对象称为\textbf{微分形式}, 由 Élie Cartan 于 1899 年引入. 它们具有自然的代数运算和微分运算, 因而是微分几何的一项基本工具.

微分形式除了可用于积分, 还出现在微分几何和物理学中的许多其他地方. 例如, 它们可以很高效地用来计算 Riemann 流形 (第~\RGref{chap:curvature}~章) 或 Lorentz 流形 (第~\RGref{chap:relativity}~章) 的曲率, 表述 Hamilton 力学 (第~\RGref{chap:geometric-mechanics}~章), 或简洁优雅地写出电磁学的 Maxwell 方程.

第~\RGref{sec:tensors}~节建立微分形式的代数结构, 回顾\textbf{张量}和\textbf{张量积}的概念, 并引入\textbf{交错张量}及其\textbf{外积}.

第~\RGref{sec:tensor-fields}~节讨论\textbf{张量场}, 它是向量场的自然推广. 这一节定义一个新运算: 协变张量场在光滑映射下的\textbf{拉回}. 第~\RGref{sec:differential-forms}~节把微分形式引入为交错张量场, 并介绍其\textbf{外微分}. 本书余下部分不以 Poincaré 引理, de Rham 上同调或 Lie 导数等重要思想为中心, 因而在习题中讨论它们.

第~\RGref{sec:integration-on-manifolds}~节定义光滑流形上的微分形式的\textbf{积分}. 这利用了微分几何的另一项基本工具, 即\textbf{单位分解}的存在性.

第~\RGref{sec:stokes-theorem}~节证明著名的 \textbf{Stokes 定理}. 它推广了向量微积分的基本定理 (Green 定理, 散度定理和向量场的经典 Stokes 定理). 它的一些推论, 如闭形式积分的\textbf{同伦不变性}, 或 \textbf{Brouwer 不动点定理}, 将在习题中探讨.

最后, 第~\RGref{sec:orientation-and-volume-forms}~节研究可定向性与一种称为\textbf{体积形式}的特殊微分形式的存在性之间的关系. 体积形式可以用来在可定向流形上定义体积概念.
